% !TEX root = ../main.tex
\section{Hallucination as Ungrounded Assertion}
\label{sec:taxonomy}

The architecture of Section~\ref{sec:architecture} produces failures with a specific
shape: assertion is computed over surface form, with no evidence gate and no resolved
context, so output can be fluent and detached from anything at once. The received
vocabulary for those failures misdescribes them, and the misdescription has consequences
for what gets measured.

The term \emph{hallucination} is in universal use and is used in at least two
incompatible ways. In popular and much industry usage it is truth-conditional: a
hallucination is a false or fabricated output. In the natural-language-generation
literature it generally is not. There, hallucination means content unsupported by
the source, whether or not that content happens to be true, and the possibility of a
factually correct hallucination is explicitly recognised
\citep{maynez2020,ji2023}.

This paper adopts the second convention and claims no priority for it. Separating
support from truth is established work, and a definition that merely restated it
would contribute nothing. The definition below departs from \emph{both} conventions
somewhere else, and it is worth being exact about where, because that departure is
this paper's first contribution and it is narrower and more specific than
``hallucination is not falsity.''

The source-based definition takes the source as given. It asks whether the assertion
follows from the evidence supplied. It does not ask the prior question: whether that
evidence \emph{governs the decision at hand}. A clinical guideline may be authentic,
authoritative, correctly retrieved, accurately quoted, and genuinely entailing of the
assertion---and concern a different jurisdiction, a superseded period, or another
patient population. Every faithfulness metric scores that assertion as clean, because
the entailment holds. Definition~\ref{def:grounding} does not, because it requires
support to be \emph{applicable at a resolved context}.

That is the move: hallucination indexed to a decision context rather than to a
source. It is the first of this paper's two central claims. The second, developed in
Section~\ref{sec:canonical-form}, is that the same architecture makes accuracy
uncertifiable. They are connected, and Section~\ref{sec:unity} states how.

\subsection{Scope of the term}
\label{sec:hall-scope}

One restriction has to be stated before the definitions, because the definition below is
broader than ordinary usage and would otherwise convict a great deal that nobody thinks
is defective. Under a requirement of applicable, traceable, dependence-bearing evidence,
the following would all count as ungrounded: a statement of common knowledge; an
inference a reader would regard as obvious; a fact the user supplied and the system
restated; hypothetical or exploratory content; and a true claim from parametric memory
in a setting where nobody expected an evidence record. Calling those hallucinations
would be a misuse of the word.

So the term is used here in an operationally restricted sense: \textbf{an asserted
proposition, presented as decision-relevant, that lacks an auditable and applicable
evidentiary basis.} The restriction is carried by \emph{presented as
decision-relevant}, and it does two things. It excludes the benign cases above, which
are not offered as grounds for consequential action. And it makes the reliance tier of
Section~\ref{sec:tiers} part of the definition's scope rather than an afterthought: the
same sentence can be unobjectionable at Tier~1 and a hallucination at Tier~3, because
what changed is what it is being relied upon for.

Readers who find the appropriation of the word uncomfortable may substitute
\emph{unwarranted assertion} or \emph{reliance-ineligible assertion} throughout without
affecting any argument. I keep ``hallucination'' because it is the word the field,
industry, and regulators actually use, and a redefinition that abandons the term
forfeits the conversation it is trying to correct. Nothing here is intended as a general
taxonomy of undesirable model output.

\subsection{Assertion}

I will talk about the deployed system $S$ rather than the model $M$ in isolation.
That is the thing being validated, and it is the thing whose failures have owners.

\begin{definition}[Assertion]
\label{def:assertion}
$\Out_S(p)$ holds iff system $S$, under input $x$, emits $p$ in an unqualified,
decision-relevant form under a stated output-interpretation policy.
\end{definition}

The definition is behavioural rather than probabilistic, and the earlier temptation to
define it by a sequence-probability threshold should be resisted for three reasons. The
probability of a full proposition's token sequence is small even when the system plainly
asserts it; equivalent phrasings of one assertion carry different sequence probabilities,
which reintroduces the form dependence of Section~\ref{sec:canonical-form} into the
definition of assertion itself; and most deployed interfaces do not expose the relevant
distribution. What makes an output an assertion is that it was emitted unqualified in a
decision-relevant position, which an interpretation policy can specify and a reviewer can
apply.

The policy has to separate commitment from mention: a system that lists $p$ among
possibilities, hedges it, or attributes it to a source without adopting it has not asserted
it. Where a graded notion of commitment is wanted it should be introduced separately
rather than folded into $\Out_S$. Nothing here requires a view about what
$M$ believes, per Section~\ref{sec:premise}, and every predicate below is likewise
defined over instrumentable behaviour.

\subsection{Grounding}
\label{sec:grounding-def}

Grounding is the property that carries the paper, and it must mean considerably
more than that a document was retrieved or a citation was emitted.

\begin{definition}[Grounded evidence]
\label{def:grounded-evidence}
For an evidential set $E$, proposition $p$, and resolved context $c$:
\begin{equation}
  \begin{aligned}
  \Grounded(E, p, c) \;\leftrightarrow\;
    &\Authentic(E) \wedge \Versioned(E) \wedge \Authoritative(E, c)
      \wedge \Applicable(E, c)\\
    &\wedge\ \Supports(E, p, c) \wedge \Sufficient(E, p, c)
      \wedge \Traceable(E, p, c).
  \end{aligned}
  \label{eq:grounded}
\end{equation}
\end{definition}

\begin{definition}[Grounded assertion]
\label{def:grounding}
$\Grd_S(p \mid c)$ holds iff there is some $E \subseteq E_S$ in the system's
evidential base such that
\begin{equation}
  \Grounded(E, p, c) \ \wedge\ \Sensitive(S, p, E) \ \wedge\ \NoDefeater(E_S, p, c),
  \label{eq:grd}
\end{equation}
where $\Sensitive(S, p, E)$ holds iff, had $E$ not supported $p$, the assertion would
not have been produced, and
\begin{equation}
  \NoDefeater(E, E_S, p, c) \;\equiv\; \DefeaterSearch(E, E_S, p, c) = \emptyset .
  \label{eq:nodefeater}
\end{equation}
$\DefeaterSearch$ is a \emph{specified search} over the governed base, not a
quantification over its subsets. It returns evidence $E'$ that is applicable at $c$, of
standing at least equal to that of $E$, and that either supports $\neg p$ or undermines
$E$'s support for $p$:
\begin{multline}
  \DefeaterSearch(E, E_S, p, c) = \bigl\{\, E' \in \mathcal{Q}(E_S, c) \;:\;
    \standing(E', c) \succeq \standing(E, c)\\
    \wedge\ \bigl[\Supports(E', \neg p, c)
      \ \vee\ \text{undermines}(E', E, p, c)\bigr] \,\bigr\}.
  \label{eq:defeatersearch}
\end{multline}
\end{definition}

The seven conjuncts of \eqref{eq:grounded} are properties of the \emph{evidence};
$\Sensitive$ is a property of the \emph{system}; and $\NoDefeater$ is a property of
the \emph{evidential base as a whole}. All three kinds are required, and the
separation matters operationally because they are checked by different means and fail
for different reasons. Section~\ref{sec:grounding-context} takes the evidence
dimensions one at a time; here I note why the other two cannot be dropped.

\subsection{Why grounding needs a defeater condition}
\label{sec:why-defeater}

$\NoDefeater$ closes a hole that the existential quantifier in \eqref{eq:grd} would
otherwise open. Without it, grounding is satisfied whenever \emph{some} applicable
authoritative source supports $p$---even if another source of equal standing in the same
base supports $\neg p$. That is not a corner case. Two regulators may differ; a standard
and its national implementation may differ; a clinical guideline and a local protocol
may differ; a group policy and a subsidiary policy may differ. A system that retrieves
one side of such a disagreement and asserts it has not produced a grounded assertion. It
has produced an arbitrary selection from a contested record.

Three features of \eqref{eq:nodefeater} and \eqref{eq:defeatersearch} are deliberate,
and each answers an objection the unbounded form invites.

\paragraph{It is a search, not a quantifier.} $\mathcal{Q}(E_S, c)$ is the governed
retrieval over the base under the applicability predicates of $c$---the same retrieval
the system already performs, run in the contrary direction. Quantifying over subsets of
a corpus would be neither computable nor useful: in a large regulatory base something
can be read as supporting almost any negation. Bounding the condition to a specified
search makes it a control that a system can execute and a reviewer can re-execute, and
it aligns the definition with the architecture of
Section~\ref{sec:governed-rag}. The cost is honest and worth stating: a defeater outside
the search is a corpus-coverage failure rather than a loophole in the definition, and it
is detected the way coverage failures are always detected---by widening the search, not
by strengthening the logic.

\paragraph{What defeats is \emph{unresolved} defeat.} Institutions act amid
disagreement constantly, and they do it legitimately, through precedence rules and
escalation procedures that were established precisely because authorities conflict. A
condition that made any equal-standing disagreement fatal to grounded assertion would
forbid most of what regulated organisations lawfully do. So $\NoDefeater$ should be read
as: no defeater that the institution's own precedence and escalation policy leaves
unresolved. Where the policy determines an outcome, the conflict has been resolved and
the assertion can be grounded on the surviving evidence, with the resolution recorded.
Where the policy is silent or the authorities are incomparable under it, the assertion is
\textsc{Contested} and escalation is the correct output. The condition is therefore
relative to a governed artefact, not absolute---and an institution with no precedence
policy has not discovered that grounded assertion is impossible; it has discovered a
missing deliverable.

\paragraph{Standing is an ordering, not a predicate.} $\standing(E, c)$ places evidence
in a partial order of authority that the institution defines: statute above guidance,
regulator above trade body, group policy above local practice except where local is
stricter. Contrary evidence of \emph{lower} standing is not a defeater---it is resolved
by precedence, which is what precedence is for. Only contrary evidence of equal or
incomparable standing defeats. This matters practically: the unordered version would
escalate every routine conflict, and a control that escalates routine cases is switched
off. It also locates a deliverable: the precedence order is a governed artefact owned by
the policy function, and its absence is why conflicts currently resolve by retrieval
order.

\paragraph{Defeat comes in two kinds.} The vocabulary is borrowed from work on
defeasible reasoning, where the distinction is standard. A \emph{rebutting} defeater
supports the contrary claim. An \emph{undercutting} defeater leaves the contrary claim
untouched and attacks the support relation instead: the instrument was withdrawn, its
scope was narrowed by a later provision, the passage was superseded in a revision the
retrieval did not reach, or the extraction misread a table. Undercutting defeat is
common in a governed corpus and invisible to a search for contradictions, because
nothing contradicts anything---the support simply no longer holds. Both belong in
\eqref{eq:defeatersearch}.

Where $\DefeaterSearch$ returns a result, the correct outcome is adjudication or
escalation rather than assertion. Section~\ref{sec:contested} treats this as a category
in its own right.

Every output of $M$ is caused by training, so causal provenance cannot be the
criterion---it would be satisfied by everything. And $\Grounded(E, p, c)$ alone is
satisfied whenever suitable evidence happens to exist somewhere in $E_S$, whether
or not the assertion had anything to do with it. A system that retrieves an
authoritative passage and then answers from parametric memory, coincidentally
agreeing with it, satisfies \eqref{eq:grounded} and is not grounded in any sense
an auditor should accept. $\Sensitive$ is what excludes this, and its operational
test is ablation: remove or contradict the putatively supporting evidence and
observe whether the assertion survives.

$\Sufficient$ closes a different hole, and one a data-quality reviewer will find
first. Evidence can be authentic, versioned, authoritative, applicable, and genuinely
supporting, and still address only part of what the claim requires. A rule imposing
five conditions is not satisfied by a source that speaks to three of them, and an
assertion resting on such a source is grounded in every dimension except adequacy.
Completeness is a standard data-quality dimension and its omission here would be
conspicuous.

The two conjuncts are independent rather than redundant, and the reason should be
stated because the redundancy objection is otherwise natural. $\Supports$ is a
\emph{positive} relation admitting degrees: the evidence favours $p$, at some strength,
as Section~\ref{sec:entails-supports} sets out. $\Sufficient$ is a \emph{coverage}
relation: the evidence addresses every condition the claim requires. Evidence can favour
a partial claim strongly---entail it, even---while covering three of five conditions, so
neither conjunct implies the other. Their remedies differ accordingly: inadequate support
is fixed by retrieving \emph{better}, insufficient coverage by retrieving \emph{more}.

\subsection{Five categories}

Let $A_i \in \{0,1\}$ record whether assertion $i$ is true under its resolved
context $c_i$, and $G_i \in \{0,1\}$ whether it is grounded in the sense of
Definition~\ref{def:grounding}. Crossing the two gives four categories.

\begin{table}[ht]
  \centering
  \begin{tabular}{@{}cc>{\raggedright\arraybackslash}p{2.4cm}>{\raggedright\arraybackslash}p{7.4cm}@{}}
    \toprule
    $A_i$ & $G_i$ & Category & Meaning \\
    \midrule
    1 & 1 & \textsc{Sound}
      & True assertion with authenticated, applicable, traceable evidentiary
        support \\
    \addlinespace
    1 & 0 & \textsc{Lucky}
      & True assertion lacking a sufficient auditable evidentiary basis \\
    \addlinespace
    0 & 1 & \textsc{Error}
      & Grounded but false: the system accurately reported what applicable
        authoritative evidence says, and the evidence diverges from the world \\
    \addlinespace
    0 & 0 & \textsc{Fabrication}
      & False assertion without an auditable evidentiary basis \\
    \bottomrule
  \end{tabular}
  \caption{The four categories. Reference-based evaluation observes $A_i$ only.}
  \label{tab:taxonomy}
\end{table}

Because grounding requires applicability, a fifth category is analytically
distinct and, in regulated deployment, the most consequential.

\begin{table}[ht]
  \centering
  \begin{tabular}{@{}>{\raggedright\arraybackslash}p{5.6cm}>{\raggedright\arraybackslash}p{2.2cm}>{\raggedright\arraybackslash}p{5.4cm}@{}}
    \toprule
    Condition & Category & Meaning \\
    \midrule
    $\Authentic(E) \wedge \Authoritative(E, c') \wedge \neg\,\Applicable(E, c)$
      & \textsc{Misapplied}
      & Authentic and perhaps authoritative evidence that does not govern $c$:
        wrong jurisdiction, time period, population, policy, product, or decision
        use \\
    \bottomrule
  \end{tabular}
  \caption{The fifth category. Operationally this is a subset of
    $G_i = 0$, since Definition~\ref{def:grounding} requires applicability. It is
    kept separate because it identifies a different control failure with a
    different owner.}
  \label{tab:misapplied}
\end{table}

A sixth follows from $\NoDefeater$ for the same reason:

\begin{itemize}
  \item[\textsc{(vi)}] \textsc{Contested} --- the governed defeater search returns
        applicable evidence of equal or incomparable standing on the other side. No
        selection is grounded, whichever side the system asserts and whichever turns out
        to be right.
\end{itemize}

Like \textsc{Misapplied}, this is operationally a subset of $G_i = 0$: the no-defeater
condition of Definition~\ref{def:grounding} fails, so the assertion is ungrounded and
falls in the right-hand column of Table~\ref{tab:taxonomy}. It is kept separate for the
same reason---the control that fails is different, the remedy is different, and the
owner is different. There is no double counting in either case.

The relationship between Table~\ref{tab:misapplied} and Table~\ref{tab:taxonomy}
should be stated explicitly, since it is a place where a reader may reasonably
suspect double-counting. Grounding as defined requires applicability, so a
misapplied assertion is \emph{operationally ungrounded}: $G_i = 0$, and it falls in
the right-hand column. It remains analytically distinct because the underlying
control failure is different in kind. In cells (ii) and (iv) no evidence was
brought to bear at all. In the misapplied case authentic evidence was retrieved,
cited, and genuinely relied upon---and its scope did not match the assertion's
resolved context. The first is a generation failure. The second is a binding
failure. Work on the generation path will rarely find it, because the information that
would reveal the mismatch---which regime governs, which version was in force at
$t_v$---is metadata about the evidence rather than content of the assertion, and is
usually not in the context at all.

\subsection{The definition}
\label{sec:hall-def}

\begin{definition}[Hallucination, error, misapplied evidence]
\label{def:hallucination}
\begin{align}
  \Hall(p \mid c)       &\;=_{\mathrm{def}}\; \Out_S(p) \wedge \neg\,\Grd_S(p \mid c), \label{eq:hall}\\
  \Err(p \mid c)        &\;=_{\mathrm{def}}\; \Out_S(p) \wedge \Grd_S(p \mid c)
                           \wedge \neg\,\True(p \mid c), \label{eq:err}\\
  \Misapplied(p \mid c) &\;=_{\mathrm{def}}\; \Out_S(p) \wedge \Supports(E, p, c')
                           \wedge \Authentic(E) \wedge \neg\,\Applicable(E, c),
                           \label{eq:misapplied}\\
  \Contested(p \mid c)  &\;=_{\mathrm{def}}\; \Out_S(p) \wedge \neg\,\NoDefeater(E, E_S, p, c).
                           \label{eq:contested}
\end{align}
\end{definition}

Note what \eqref{eq:hall} omits: \textbf{any condition on the truth of $p$.} A
hallucination is an assertion unsupported by applicable evidence, and that is all.
Truth and falsity distinguish successes from errors; they do not distinguish
hallucination from sound assertion.

Two consequences follow immediately, and both invert the standard usage. A false
assertion faithfully reporting an applicable source is \emph{not} a hallucination;
it is an error, and the fault lies in the evidence. A true assertion resting on
nothing \emph{is} a hallucination, and remains one however often it turns out
right. The next three subsections defend each of these, and the defence is in each
case an argument about remedies rather than about words.

The definition also has a property that neither convention has: it is
\emph{evaluable at the point of assertion}, without an answer key. That is the
subject of Section~\ref{sec:auditable}, and it is why this is an engineering claim
rather than a terminological preference. A truth-conditional definition can only be
applied offline, against known answers. A source-based definition can be applied at
runtime but only relative to whatever was retrieved. Definition~\ref{def:grounding}
is applied at runtime against the evidence \emph{and} the context, which is what a
gate needs.

\subsection{What is claimed as new, and what is not}
\label{sec:novelty}

Groundedness, provenance, factuality, attribution, and contextual relevance all have
substantial literatures, and nothing here is offered as a discovery in any of them.
The claim is narrower, and stating it precisely is worth a paragraph because a claim
pitched wider would be false.

This paper does not treat hallucination as synonymous with factual error, and does
not claim priority for that separation. It formalises the relevant failure as the
absence of an auditable relation between an assertion, an authenticated and
\emph{contextually applicable} evidentiary basis, and a verification pathway that an
institution has authorised. Four components of that formalisation are, so far as I am
aware, not present in existing definitions:

\begin{enumerate}
  \item \textbf{Applicability as a conjunct.} Support is indexed to
        $c = \{j, t_v, f, u\}$, so evidence that entails the assertion but does not
        govern the case fails the definition. Faithfulness metrics take the source as
        given and cannot express this condition.
  \item \textbf{Dependence rather than entailment.} $\Sensitive$ requires the
        assertion to have been produced \emph{because of} the evidence, not merely to
        be consistent with it. A system that retrieves an authoritative passage and
        answers from parametric memory satisfies entailment and fails this.
  \item \textbf{Authority and version as conditions on support.} Whether a source has
        standing for this decision class, and which version of it applied at $t$, are
        treated as preconditions of valid support rather than as metadata.
  \item \textbf{Defeat by contrary authority.} $\NoDefeater$ makes grounding a
        property of the evidential base rather than of a selected subset of it, so an
        assertion drawn from one side of a live disagreement between applicable
        authorities is not grounded. Faithfulness and attribution metrics evaluate an
        assertion against the source it cites and have no way to express this.
  \item \textbf{The non-identification result.} Proposition~\ref{prop:independence}
        establishes that no accuracy figure constrains the resulting rate, which is
        what makes the redefinition consequential for validation rather than
        expository.
\end{enumerate}

Taken together these separate four cases that pooled accuracy and conventional
source-faithfulness metrics conflate: lucky truth, grounded error, fabrication, and
contextual misapplication. The compact form of the claim is that factual accuracy is
insufficient for institutional reliance, because an assertion must also possess a
reconstructable and contextually applicable evidentiary lineage---and that a metric
which cannot distinguish those four cases cannot support a reliance decision.

% TODO(priority): I have not been able to establish exhaustively that no existing
% work indexes a grounding or hallucination definition to a decision context in this
% way. Before asserting priority in items 1-3, run a targeted check of the
% temporal-QA, legal-NLP, and clinical-decision-support literatures, where the
% individual phenomena (temporal validity, jurisdiction, population applicability)
% are certainly recognised even if their integration into a definition is not.

\subsection{Cell (iii): a faithful report of a bad source is an error, not a fabrication}
\label{sec:cell-iii}

A corpus contains, consistently and across several reference works, an incorrect
effective date for a regulatory provision. The system reports that date. Ablation
confirms the assertion is driven by those documents: remove them and it stops.

Reference-based evaluation scores this as a hallucination, because the output is
false.
That classification discards a distinction the organisation needs. The system did
what a reporting component is supposed to do: it reported what its evidence said,
through a mechanism sensitive to that evidence. The defect is in the evidence.

One clarification is owed first, because the cell looks over-full as often stated. If
grounding requires $\Supports$ and $\Sufficient$, how can a grounded assertion be false?
Several candidate causes turn out not to belong here. A misread passage means the evidence
does not in fact support the claim, so $\Supports$ fails and the assertion is ungrounded.
An invalid inference is the same. Overstating a source fails $\Sufficient$. Those are
\textsc{Fabrication} in the sense of \eqref{eq:hall}, not \textsc{Error}, and their remedy
is on the generation path.

\textsc{Error} is therefore narrower than it first appears, and the narrowness is the
point: it is the case where the system accurately reported applicable authoritative
evidence and \emph{the evidence itself diverges from the world}. Grounding, as defined, is
institutional evidentiary grounding rather than truth-tracking, so the cell is exactly the
gap between the two. It is the only failure the model and the pipeline cannot be blamed
for, and the only one that no amount of generation-path work will reduce.

The distinction is not terminological, it is budgetary. Cell (iii) is addressed by
work on the evidential base---source selection, authority ranking, curation,
retrieval from systems of record rather than secondary literature. Cell (iv) is
addressed by work on the generation path---evidence gating, abstention policies,
refusal thresholds. Different teams, different skills, different funding cycles, as
Section~\ref{sec:three-teams} develops. A metric that pools them reports a number
on which no one can act.

There is a normative point as well. We do not say that a careful analyst
fabricated when she accurately reported a source that later proved wrong. We say
she was misled, and we look at her sources. Machine assertion should be held to the
same distinction, in both directions.

\subsection{Cell (ii): a lucky right answer is not evidence about the system}
\label{sec:cell-ii}

Now the reverse. The system is asked for a figure absent from its evidential
base---the headcount of a small subsidiary. The generation path, running on surface
regularities about how entities of this type and sentences of this genre continue,
produces $8{,}200$. That happens to be right.

Reference-based evaluation scores this as a success. Every property that makes
fabrication dangerous is nonetheless present, and the precise defect is worth
stating because it is what makes the case more than a curiosity: had the true
figure been $3{,}400$, the same process would have produced the same output. The
assertion does not depend on the fact it reports. It is not a measurement; it is a
draw from a distribution that happened to land on the answer.

The consequence for validation is the one that matters. A user who receives a
cell-(ii) output, checks it, and finds it correct has acquired \emph{no} evidence
about the system. The check confirmed the output, not the process, and the process
will behave identically on the next question, where the draw may land elsewhere.
Lucky answers are not repeatable performances. Grounding is---and that asymmetry
is the argument of Section~\ref{sec:auditable}, arriving early.

This also explains why a benchmark that rewards cell (ii) does more than fail to
inform. Scoring lucky answers as successes applies
optimisation pressure toward confident guessing and away from abstention, which by
Section~\ref{sec:rlhf} is a direction the training regime already favours. The
metric compounds a bias the architecture supplies.

\subsection{Cell (v): authentic is not applicable}
\label{sec:cell-v}

The category pooled metrics hide, and in my experience the most common serious
failure in regulated deployment.

The system asserts that a treatment is indicated for a condition. The assertion is
driven by clinical guidance that says exactly that, and ablation confirms it. But
the guidance concerns an adult population and the case is paediatric. Or a system
asserts a capital treatment on the basis of a rule that governs a different
jurisdiction. Or it cites a policy that was superseded before the decision date. Or
it applies a product rule to a different product class.

Without the context index these are cell (iii): grounded, false, therefore charged
to a defective evidential base. That verdict is wrong twice over. The base was not
defective---the guidance is accurate for its population, the foreign rule is
correctly stated, the superseded policy was valid when issued. And the implied
remedy is wrong: curating the corpus will not help, because nothing in it is
incorrect.

What failed is applicability, and the remedy is context resolution: establishing
$c$ before answering, checking currency against $t$, checking jurisdiction,
material facts, and intended use, and abstaining or qualifying when the applicable
evidence has not been identified. These examples are illustrations of the failure
shape rather than reports of measured incidents, and
Section~\ref{sec:grounding-context} develops the underlying generalisation error.
\subsection{Cell (vi): a contested record is not a grounded one}
\label{sec:contested}

The failure that the existential quantifier hides, and the one most likely to be
misclassified as an error.

Two authoritative sources in the base disagree. A group accounting policy and a
statutory requirement in one jurisdiction reach different treatments. A clinical
guideline and a local formulary differ on a first-line therapy. A standard and its
national implementation diverge on a threshold. The system retrieves one, and asserts
it. Ablation confirms the assertion rests on that source; the source is authentic,
versioned, authoritative, applicable, sufficient, and traceable.

Every dimension of \eqref{eq:grounded} passes, and the assertion is not fit for
reliance. What the system has done is resolve a policy conflict by retrieval order,
which is to say arbitrarily, and it has produced a record showing a single clean
citation with no indication that the question was contested at all. That is worse than
an unsupported answer, because the audit trail actively conceals the problem: a
reviewer sees one authoritative source and an assertion that follows from it.

Note that the defect is independent of which side is correct. If the retrieved source
happens to state the treatment a reviewer would have chosen, nothing has been
established---the selection was still arbitrary, and the next case may go the other
way. This is the lucky cell again, at the level of source selection rather than token
generation.

The remedy is neither corpus work nor generation work. It is \emph{adjudication}: a
rule that determines which authority governs when authorities conflict, and an
escalation path for cases the rule does not settle. That is a policy artefact, and its
absence is why contested cases surface as apparently well-grounded assertions.

The categories exist to route work, and this is what they are for. Four remedies,
across three owners:
\label{sec:remedies}

\begin{table}[ht]
  \centering
  \begin{tabular}{@{}>{\raggedright\arraybackslash}p{2.6cm}>{\raggedright\arraybackslash}p{4.0cm}>{\raggedright\arraybackslash}p{4.4cm}>{\raggedright\arraybackslash}p{2.3cm}@{}}
    \toprule
    Category & Diagnosis & Remedy & Typical owner \\
    \midrule
    \textsc{Error}
      & applicable evidence, but wrong
      & source authority, curation, systems of record
      & data \\
    \addlinespace
    \textsc{Misapplied}
      & evidence inapplicable at $c$
      & context resolution, currency and scope binding
      & policy / risk \\
    \addlinespace
    \textsc{Contested}
      & applicable authority on both sides
      & precedence rules, adjudication, escalation path
      & policy / risk \\
    \addlinespace
    \textsc{Lucky}, \textsc{Fabrication}
      & no evidential connection
      & evidence gating, abstention, refusal thresholds
      & platform \\
    \bottomrule
  \end{tabular}
  \caption{A single pooled hallucination rate cannot distinguish these rows, and
    therefore cannot tell an organisation which of four budgets to spend.
    Section~\ref{sec:three-teams} argues that the division of ownership is itself
    a principal source of failure.}
  \label{tab:remedies}
\end{table}
